\documentclass[12pt]{article}
\usepackage[legalpaper,
            lmargin=1in,rmargin=0.5in,
            bmargin=0.5in,tmargin=0.8in]{geometry}
\usepackage{fontspec}
\usepackage{titlesec}
\usepackage{multirow}
\usepackage[colorlinks=true,urlcolor=Blue]{hyperref}
\usepackage{graphicx}
\usepackage{array}
\usepackage{makecell}
\usepackage{fancyhdr}
\usepackage[none]{hyphenat}
\usepackage{longtable}
\usepackage[dvipsnames]{xcolor}
\usepackage[banglamainfont=Kalpurush,
            banglattfont=SolaimanLipi,
            % feature=1,
            % changecounternumbering=0
           ]{latexbangla}

\let\oldhref\href
\renewcommand{\href}[3][black]{\oldhref{#2}{\color{#1}{#3}}}
\newcommand{\my}{\includegraphics[height=3.2em]{pic/my.png}}

\pagestyle{fancy}
\fancyhf{}
\renewcommand{\headrulewidth}{0pt}
% header
\chead{
\underline{{পৃষ্ঠা - \thepage}}
\\
}
\rhead{{\filenou}}
% footer
% \rfoot{চলমান পৃষ্ঠা-\thepage}

\fancypagestyle{slogan}
{
\fancyhf{}
\renewcommand{\headrulewidth}{0pt}
% header
\lhead{
\framebox[1.1\width]{\footnotesize{``উন্নয়নের অক্সিজেন রাজস্ব''}}
}
\rhead{
\my
\\
\framebox[1.1\width]{\footnotesize{``জনকল্যানে রাজস্ব''}}
}
}


% \AtEndDocument{\thispagestyle{empty}}
\newcommand{\fileno}{নথি নং - ২৬৫৭/এপি/সেকশন-৪/২০২১-২০২২}
\newcommand{\filenou}{\underline{\fileno}}
\newcommand{\fdt}{\hspace*{3em} তারিখ: \hspace{2.4em} /১২/২১}
\newcommand{\rodt}{তারিখ: \hspace{2.0em}/১২/২০২১ ইং}
\newcommand{\product}{}
\newcommand{\good}{VISCOSE RAYON STAPLE FIBER}
\newcommand{\pkg}{836.00 BALES=258,398.80 KGS}
\newcommand{\co}{INDONESIA}
\newcommand{\coship}{INDONESIA}
\newcommand{\vessel}{CAPE SYROS}
\newcommand{\blno}{SINA153435}
\newcommand{\bldt}{08.11.2021}
\newcommand{\beno}{C-1964404}
\newcommand{\bedt}{06.12.21}
\newcommand{\lcno}{0000000121230397}
\newcommand{\lcdt}{13.10.2021}
\newcommand{\lcano}{128404}
\newcommand{\lcadt}{\lcdt}
\newcommand{\lienbank}{AGRANI BANK LTD.}
\newcommand{\invno}{5960002158}
\newcommand{\invdt}{08.11.2021}
\newcommand{\ssml}{SHAMEEM SPINNING MILLS LTD.}
\newcommand{\ssmla}{SHAFIPUR, KALIAKAIR
\\
GAZIPUR-1751, BANGLADESH.}
\newcommand{\jsml}{JAMUNA SPINNING MILLS LTD.}
\newcommand{\jsmlau}{JAMUNA SPINNING MILLS LTD. U-2
\\
BEJURA, SOUTH BEJURA,
\\
MADHABPUR, HOBIGONJ-3331
}
\newcommand{\jsmlausl}{JAMUNA SPINNING MILLS LTD. U-2,
BEJURA, SOUTH BEJURA,
MADHABPUR, HOBIGONJ-3331}
\newcommand{\jdl}{JAMUNA DENIMS LTD.}
\newcommand{\jdla}{SHAFIPUR, KALIAKAIR
\\
GAZIPUR-1751, BANGLADESH.}
\newcommand{\scml}{SHAMEEM COMPOSITE MILLS LTD.}
\newcommand{\jeal}{JAMUNA ELECTRONICS \& AUTOMOBILES LTD.}
\newcommand{\jdwl}{JAMUNA DENIMS \& WEAVING LTD.}
\newcommand{\htf}{HOORAIN HTF LTD.}
\newcommand{\tdj}{M/S THE DAILY JUGANTOR}
\newcommand{\tdja}{Ka-244, Kuril, Progati Sarani, Vatara PS
\\
Dhaka-1229,Bangladesh}

\newcommand{\impn}{\jsml}
\newcommand{\impadd}{\jsmlau}
\newcommand{\impaddsl}{\jsmlausl}
\newcommand{\impbin}{0001239819-0704}
\newcommand{\cnfn}{SHAMEEM SPINNING MILLS LTD.}
\newcommand{\cnfadd}{92, HIGH LEVEL ROAD
\newline
LALKHAN BAZAR, CHITTAGONG}
\newcommand{\cnfain}{301 08 3417}
\newcommand{\crf}{NON CRF}
\newcommand{\crfdt}{}
\newcommand{\ircno}{আইআরসি নং - ২৬০৩২৬১২০৪২৬৭২০}
\newcommand{\ircrenewdt}{৩০/০৬/২০২২ ইং}
\newcommand{\musokr}{সেপ্টেম্বর-২১}
\newcommand{\hscode}{55041000}
\newcommand{\price}{US\$ 486,705.96}
\newcommand{\menifest}{2021/5594}
\newcommand{\srooot}{এসআরও নং -১২০-আইন/২০২১/০৯/কাস্টমস}
\newcommand{\sroootd}{তারিখ: ২৪/০৫/২০২১ খ্রি:}
\newcommand{\srooof}{এসআরও নং -১১৪-আইন/২০২১/০৩/কাস্টমস}
\newcommand{\srooofd}{তারিখ: ২৪/০৫/২০২১ খ্রি:}
\newcommand{\srotsz}{এসআরও নং - ৩৬০-আইন/২০১৩/২৪৬৮/কাস্টমস}
\newcommand{\srotszd}{তারিখ: ২৫/১১/২০১৩ খ্রি:}
\newcommand{\srofs}{এসআরও নং - ৫৭-আইন/২০০০/১৮২১/শুল্ক}
\newcommand{\srofsd}{তারিখ: ২৩/০২/২০০০ খ্রি:}
\newcommand{\nbrosnt}{জাতীয় রাজস্ব বোর্ডের পত্র নং - ০৮.০১.০০০০.০৬৮.১৮.০০৪.১৭/১৬৯(৩)}
\newcommand{\nbrosntd}{তারিখ: ২৭/০৬/২০২১ খ্রি:}
\newcommand{\nbrfs}{জাতীয় রাজস্ব বোর্ডের পত্র নং - ০৮.০১.০০০০.০৩৪.০২.৩০১.১৯-৪৭}
\newcommand{\nbrfsd}{তারিখ: ১২/০৭/২০২১ খ্রি:}
\newcommand{\btmaltno}{প্রত্যয়ন পত্র নং- ২৮৫১১}
\newcommand{\btmaltnodt}{তারিখ: ১৬/১১/২০২১ খ্রি:}
\newcommand{\cpc}{(সিপিসি ৪০০০/৪০৯)}
\newcommand{\allotment}{৭৫০০}
\newcommand{\relesed}{১৭৫.৬৩৯}
\newcommand{\current}{২৫৮.৩৯৮৮}
\newcommand{\remaining}{৭৩২৪.৩৬}
\newcommand{\balance}{৭০৬৫.৯৬১২}
\newcommand{\taxtab}{
\begin{longtable}{|c|c|c|c|c|c|c|c|}
\hline
\textbf{
\makecell{
ক্রঃ \\ নং
}
}
&
\textbf{
\makecell{
পণ্যের বর্ণনা
}
}
&
\textbf{
\makecell{
পরিমাণ
}
}
& \textbf{
\makecell{
ইনভয়েস
ঘোষিত
\\
এইচএসকোড
\\
(শুল্কহারসহ)
}
}
&
\textbf{
\makecell{
প্রকৃত
\\
এইচএসকোড
\\
(শুল্কহারসহ)
}
}
&
\textbf{
\makecell{
ইনভয়েস
\\
প্রত্যায়িত
\\
একক মূল্য
\\
(US\$)
}
}
&
\textbf{
\makecell{
সমসাময়িক
কালের
\\
অভিন্ন/অনুরুপ
\\
পণ্যের
\\
একক মূল্য
(US\$)
}
}
&
\textbf{
\makecell{
প্রস্তাবিত
\\
একক মূল্য
\\
(US\$)
}
} \\
\hline
% row
\makecell{
01
}
&
\makecell{
VISCOSE RAYON STAPLE FIBER
\\
C/O. INDONESIA
}
&
\makecell{
258,398.80
\\
KGS
}
&
\makecell{
55041000
\\
BTMA LETTER
\\
SRO-120/2021
\\
CPC 4000/157
}
&
\makecell{
55041000
\\
BTMA LETTER
\\
SRO-120/2021
\\
CPC 4000/157
}
&
\makecell{
US\$
\\
1.88/KG
}
&
\makecell{
US\$
\\
2.55/KG
\\
DATABASE
VALUE
\\
C-1036262
\\
DT:25.06.2021
}
&
\makecell{
US\$
\\
2.55/KG
} \\
\hline
\end{longtable}
}
\begin{document}
\noindent
\begin{minipage}[t]{0.05\linewidth}
% ek
০১।
\end{minipage}
\begin{minipage}[t]{0.95\linewidth}
বি/ই রেজি: নং- {\beno}, তারিখ: {\bedt}
নথিভূক্ত করে
পরবর্তী কার্যক্রমের জন্য উপস্থাপন করা হলো।
\\
\\
\end{minipage}
\begin{minipage}[t]{0.05\linewidth}
\hspace*{0em}
\end{minipage}
\begin{minipage}[t]{0.05\linewidth}
সহকারী
\end{minipage}
\begin{minipage}[t]{0.37\linewidth}
\hspace{0em}
\end{minipage}
\begin{minipage}[t]{0.53\linewidth}
\textbf{শুল্কায়ন কর্মকর্তা}
\\
\end{minipage}
\begin{minipage}[t]{0.05\linewidth}
% dui
০২।
\end{minipage}
\begin{minipage}[t]{0.95\linewidth}
\underline{\textbf {আমদানিকৃত পণ্য চালানের
মৌলিক তথ্য:}}
\\
\end{minipage}
\footnotesize
\begin{minipage}[t]{0.05\linewidth}
\hspace*{1em}
\end{minipage}
\begin{minipage}[t]{0.40\linewidth}
(ক) বি/ই রেজি: নং ও তারিখ
\end{minipage}
\begin{minipage}[t]{0.02\linewidth}
:
\end{minipage}
\begin{minipage}[t]{0.53\linewidth}
\textbf{{\beno}} \hspace{2em} DT: {\bedt}
\\
\end{minipage}
\begin{minipage}[t]{0.05\linewidth}
\hspace*{1em}
\end{minipage}
\begin{minipage}[t]{0.40\linewidth}
(খ) আমদানিকারকের নাম, ঠিকানা
ও BIN নম্বর
\end{minipage}
\begin{minipage}[t]{0.02\linewidth}
:
\end{minipage}
\begin{minipage}[t]{0.53\linewidth}
\textbf{{\impn}}
\\
{\impadd}
\\
BIN NO. {\impbin}
\\
\end{minipage}
\begin{minipage}[t]{0.05\linewidth}
\hspace*{1em}
\end{minipage}
\begin{minipage}[t]{0.40\linewidth}
(গ) সিএন্ডএফ এজেন্টের নাম, ঠিকানা
ও AIN নম্বর
\end{minipage}
\begin{minipage}[t]{0.02\linewidth}
:
\end{minipage}
\begin{minipage}[t]{0.53\linewidth}
\textbf{{\cnfn}}
\\
{\cnfadd}
\\
AIN NO. {\cnfain}
\\
\end{minipage}
\begin{minipage}[t]{0.05\linewidth}
\hspace*{1em}
\end{minipage}
\begin{minipage}[t]{0.40\linewidth}
(ঘ) এল/সি নং ও তারিখ
\end{minipage}
\begin{minipage}[t]{0.02\linewidth}
:
\end{minipage}
\begin{minipage}[t]{0.53\linewidth}
{\lcno} \hspace{2em} DT: {\lcdt}
\\
\end{minipage}
\begin{minipage}[t]{0.05\linewidth}
\hspace*{1em}
\end{minipage}
\begin{minipage}[t]{0.40\linewidth}
(ঙ) লিয়েন ব্যাংকের নাম
\end{minipage}
\begin{minipage}[t]{0.02\linewidth}
:
\end{minipage}
\begin{minipage}[t]{0.53\linewidth}
{\lienbank}
\\
\end{minipage}
\begin{minipage}[t]{0.05\linewidth}
\hspace*{1em}
\end{minipage}
\begin{minipage}[t]{0.40\linewidth}
(চ) এলসিএ নং ও তারিখ
\end{minipage}
\begin{minipage}[t]{0.02\linewidth}
:
\end{minipage}
\begin{minipage}[t]{0.53\linewidth}
{\lcano} \hspace{2em} DT: {\lcadt}
\\
\end{minipage}
\begin{minipage}[t]{0.05\linewidth}
\hspace*{1em}
\end{minipage}
\begin{minipage}[t]{0.40\linewidth}
(ছ) বি/এল নং ও তারিখ
\end{minipage}
\begin{minipage}[t]{0.02\linewidth}
:
\end{minipage}
\begin{minipage}[t]{0.53\linewidth}
{\blno} \hspace{2em} DT: {\bldt}
\\
\end{minipage}
\begin{minipage}[t]{0.05\linewidth}
\hspace*{1em}
\end{minipage}
\begin{minipage}[t]{0.40\linewidth}
(জ) বাণিজ্যিক ইনভয়েস নং ও তারিখ
\end{minipage}
\begin{minipage}[t]{0.02\linewidth}
:
\end{minipage}
\begin{minipage}[t]{0.53\linewidth}
{\invno} \hspace{2em} DT: {\invdt}
\\
\end{minipage}
\begin{minipage}[t]{0.05\linewidth}
\hspace*{1em}
\end{minipage}
\begin{minipage}[t]{0.40\linewidth}
(ঝ) সিআরএফ নং ও ইস্যুর তারিখ
\end{minipage}
\begin{minipage}[t]{0.02\linewidth}
:
\end{minipage}
\begin{minipage}[t]{0.53\linewidth}
{\crf} \hspace{2em} {\crfdt}
\\
\end{minipage}
\begin{minipage}[t]{0.05\linewidth}
\hspace*{1em}
\end{minipage}
\begin{minipage}[t]{0.40\linewidth}
(ঞ) পণ্যের বিবরণ
\end{minipage}
\begin{minipage}[t]{0.02\linewidth}
:
\end{minipage}
\begin{minipage}[t]{0.53\linewidth}
{\good}
\\
\end{minipage}
\begin{minipage}[t]{0.05\linewidth}
\hspace*{1em}
\end{minipage}
\begin{minipage}[t]{0.40\linewidth}
(ট) পণ্যের পরিমাণ (একক সহ)
\end{minipage}
\begin{minipage}[t]{0.02\linewidth}
:
\end{minipage}
\begin{minipage}[t]{0.53\linewidth}
{\pkg}
\\
\end{minipage}
\begin{minipage}[t]{0.05\linewidth}
\hspace*{1em}
\end{minipage}
\begin{minipage}[t]{0.40\linewidth}
(ঠ) পণ্যের এইচ.এস.কোড
\end{minipage}
\begin{minipage}[t]{0.02\linewidth}
:
\end{minipage}
\begin{minipage}[t]{0.53\linewidth}
{\hscode}
\\
\end{minipage}
\begin{minipage}[t]{0.05\linewidth}
\hspace*{1em}
\end{minipage}
\begin{minipage}[t]{0.40\linewidth}
(ড) পণ্যের মূল্য (ইনভয়েস অনুযায়ী)
\end{minipage}
\begin{minipage}[t]{0.02\linewidth}
:
\end{minipage}
\begin{minipage}[t]{0.53\linewidth}
{\price}
\\
\end{minipage}
\begin{minipage}[t]{0.05\linewidth}
\hspace*{1em}
\end{minipage}
\begin{minipage}[t]{0.40\linewidth}
(ঢ) কান্ট্রি অব অরিজিন
\end{minipage}
\begin{minipage}[t]{0.02\linewidth}
:
\end{minipage}
\begin{minipage}[t]{0.53\linewidth}
{\co}
\\
\end{minipage}
\begin{minipage}[t]{0.05\linewidth}
\hspace*{1em}
\end{minipage}
\begin{minipage}[t]{0.40\linewidth}
(ণ) কান্ট্রি অব শিপমেন্ট
\end{minipage}
\begin{minipage}[t]{0.02\linewidth}
:
\end{minipage}
\begin{minipage}[t]{0.53\linewidth}
{\coship}
\\
\end{minipage}
\begin{minipage}[t]{0.05\linewidth}
\hspace*{1em}
\end{minipage}
\begin{minipage}[t]{0.40\linewidth}
(ত) জাহাজের নাম
\end{minipage}
\begin{minipage}[t]{0.02\linewidth}
:
\end{minipage}
\begin{minipage}[t]{0.53\linewidth}
{\vessel}
\end{minipage}
\begin{minipage}[t]{0.05\linewidth}
\hspace*{1em}
\end{minipage}
\begin{minipage}[t]{0.40\linewidth}
\hspace*{1.8em}পালা নং বি/এল নং
\end{minipage}
\begin{minipage}[t]{0.02\linewidth}
\hspace{1em}
\end{minipage}
\begin{minipage}[t]{0.53\linewidth}
{\menifest}, B/L {\blno}
\\
\end{minipage}
\begin{minipage}[t]{0.05\linewidth}
\hspace*{1em}
\end{minipage}
\begin{minipage}[t]{0.40\linewidth}
(থ) মেনিফিস্ট নং
\end{minipage}
\begin{minipage}[t]{0.02\linewidth}
:
\end{minipage}
\begin{minipage}[t]{0.53\linewidth}
{\menifest}
\\
\end{minipage}
\normalsize
\begin{minipage}[t]{0.05\linewidth}
% tin
০৩।
\end{minipage}
\begin{minipage}[t]{0.95\linewidth}
\underline{\textbf{শুল্কায়ন সেকশনের পর্যালোচনা:}}
\end{minipage}
\begin{minipage}[t]{0.05\linewidth}
\hspace{1em}
\end{minipage}
\begin{minipage}[t]{0.05\linewidth}
% tina
(ক)
\end{minipage}
\begin{minipage}[t]{0.90\linewidth}
\underline{\textbf{আমদানি দলিল পত্র যাচাই:}}
পণ্যচালান খালাসের জন্য নিম্নবর্ণিত দলিলাদিসহ বি/ই দাখিল করা
হয়েছে যা সংশ্লিষ্ট ব্যাংকের অথরাইজড কর্মকর্তা কর্তৃক সত্যায়িত।
\\
(১) এল.সি এবং এল.সি.এ ফরম।
\\
(২) ইনভয়েস।
\\
(৩) প্যাকিং লিস্ট।
\\
(৪) বি/এল।
\\
(৫) কান্ট্রি অব অরিজিন সনদ।
\\
দলিলাদি পর্যালোচনায় এগুলো
সঠিক পাওয়া যায়।
\\
\end{minipage}
\begin{minipage}[t]{0.05\linewidth}
\hspace{1em}
\end{minipage}
\begin{minipage}[t]{0.05\linewidth}
% tinb
(খ)
\end{minipage}
\begin{minipage}[t]{0.90\linewidth}
\underline{\textbf{আমদানি যোগ্যতা যাচাই:}}
প্রচলিত আমদানিনীতি আদেশ পর্যালোচনা করে দেখা যায় যে, পণ্যগুলি অবাধে আমদানিযোগ্য।
আলোচ্য চালানের ক্ষেত্রে আমদানিনীতি আদেশের প্রযোজ্য অন্যান্য শর্ত (কান্ট্রি অব অরিজিন, রেজিঃ
সার্টিফিকেট ইত্যাদি) প্রতিপালিত হয়েছে।
\\
\end{minipage}
\begin{minipage}[t]{0.05\linewidth}
\hspace{1em}
\end{minipage}
\begin{minipage}[t]{0.05\linewidth}
% tinc
(গ)
\end{minipage}
\begin{minipage}[t]{0.90\linewidth}
\underline{\textbf{কায়িক পরীক্ষার প্রতিবেদন পর্যালোচনা:}}
আলোচ্য পণ্যচালানটির ক্ষেত্রে প্রযোজ্য নয়।
\\
\end{minipage}
\begin{minipage}[t]{0.05\linewidth}
\hspace{1em}
\end{minipage}
\begin{minipage}[t]{0.05\linewidth}
% tind
(ঘ)
\end{minipage}
\begin{minipage}[t]{0.90\linewidth}
\underline{\textbf{রাসায়নিক পরীক্ষা সংক্রান্ত মন্তব্য:}}
আলোচ্য পণ্যচালানটির ক্ষেত্রে প্রযোজ্য নয়।
\\
\end{minipage}
\begin{minipage}[t]{0.05\linewidth}
\hspace{1em}
\end{minipage}
\begin{minipage}[t]{0.05\linewidth}
% tine
(ঙ)
\end{minipage}
\begin{minipage}[t]{0.90\linewidth}
\underline{\textbf{এইচ.এস.কোড সঠিকতা যাচাই:}}
The Customs Act, 1969 এর FIRST SCHEDULE ও
EXPLANATORY NOTES প্রচলিত এসআরও/স্থায়ী আদেশ ইত্যাদির আলোকে পরীক্ষা করে পণ্যের সঠিক HS Code নোট অনুচ্ছেদ-৪
এর ছকে উপস্থাপন করা হলো।
\\
\end{minipage}
\begin{minipage}[t]{0.05\linewidth}
\hspace{1em}
\end{minipage}
\begin{minipage}[t]{0.05\linewidth}
% ting
(চ)
\end{minipage}
\begin{minipage}[t]{0.90\linewidth}
\underline{\textbf{অভিযোগ সংক্রান্ত:}} আলোচ্য পণ্যচালানে
গোপন সংবাদ দাতা, শুল্ক গোযেন্দা বা
জাতীয় রাজস্ব বোর্ডের কিংবা অন্য দপ্তর থেকে
কোন অভিযোগ পাওয়া যায় নাই।
\\
\end{minipage}
\begin{minipage}[t]{0.05\linewidth}
\hspace{1em}
\end{minipage}
\begin{minipage}[t]{0.05\linewidth}
% tinh
(ছ)
\end{minipage}
\begin{minipage}[t]{0.90\linewidth}
\underline{\textbf{ন্যায় নির্ণয় সংক্রান্ত:}} প্রযোজ্য নয়।
\\
\end{minipage}
\begin{minipage}[t]{0.05\linewidth}
\hspace{1em}
\end{minipage}
\begin{minipage}[t]{0.05\linewidth}
% tinh
(জ)
\end{minipage}
\begin{minipage}[t]{0.05\linewidth}
\end{minipage}
\begin{minipage}[t]{0.90\linewidth}
\underline{\textbf{উল্লেখ করার মত প্রাসঙ্গিক অন্যান্য বিষয়:}}
\end{minipage}
\footnotesize
\begin{minipage}[t]{0.05\linewidth}
\hspace{1em}
\end{minipage}
\begin{minipage}[t]{0.05\linewidth}
\hspace{1em}
\end{minipage}
\begin{minipage}[t]{0.05\linewidth}
(১)
\end{minipage}
\begin{minipage}[t]{0.85\linewidth}
আলোচ্য পণ্যচালানটির ঘোষিত HS Code {\hscode}
যা {\srooot}, {\sroootd} এর শর্ত মোতাবেক
আমদানি করা হয়েছে।
\end{minipage}
\begin{minipage}[t]{0.05\linewidth}
\hspace{1em}
\end{minipage}
\begin{minipage}[t]{0.05\linewidth}
\hspace{1em}
\end{minipage}
\begin{minipage}[t]{0.05\linewidth}
(২)
\end{minipage}
\begin{minipage}[t]{0.85\linewidth}
আলোচ্য পণ্যচালানের বিপরীতে বাংলাদেশ
টেক্সটাইল মিলস এসোসিয়েশন এর সভাপতি
কর্তৃক স্বাক্ষরিত প্রত্যয়নপত্র
দাখিল করেছেন। প্রত্যয়নপত্রে আমদানিকৃত
যন্ত্রপাতি, যন্ত্রাংশ, ও উপকরণ
উক্ত প্রতিষ্ঠানের উৎপাদন প্রক্রিয়ায়
ব্যবহৃত হইবে এবং আমদানিকৃত
যন্ত্রপাতি, যন্ত্রাংশ, ও উপকরণ
{\srooot}, {\sroootd} অনুযায়ী
রেয়াতীহারে শুল্কায়নের সুপারিশ
করেন। বাংলাদেশ টেক্সটাইল মিলস
এসোসিয়েশনের দাখিলকৃত {\btmaltno}
{\btmaltnodt} নথির যোগাযোগ অংশে রক্ষিত
আছে, দয়া করে দেখা যেতে পারে।
\end{minipage}
\begin{minipage}[t]{0.05\linewidth}
\hspace{1em}
\end{minipage}
\begin{minipage}[t]{0.05\linewidth}
\hspace{1em}
\end{minipage}
\begin{minipage}[t]{0.05\linewidth}
(৩)
\end{minipage}
\begin{minipage}[t]{0.85\linewidth}
এসআরও শর্ত মোতাবেক আমদানিকারক ৩০০.০০ (তিনশত)
টাকার নন-জুডিশিয়াল স্ট্যাম্পে একখানা অঙ্গীকারনামা
দাখিল করেছেন। দাখিলকৃত অঙ্গীকারনামায় আমদানিকৃত
যন্ত্রপাতি, যন্ত্রাংশ, ও উপকরণ
উক্ত প্রতিষ্ঠানের উৎপাদন প্রক্রিয়ায়
ব্যবহার করবেন। অন্যথায় আমদানিকৃত উক্ত পণ্যের উপর
প্রযোজ্য স্বাভাবিক হারের শুল্ককর পরিশোধ ছাড়াও
শুল্ক আইন মোতাবেক শুল্ক কর্তৃপক্ষ
গৃহীত যে কোন আইনানুগ সিদ্ধান্ত মানিয়া নিতে বাধ্য থাকিবেন
মর্মে উল্লেখ করেছেন। আমদানিকারক কর্তৃক দাখিলকৃত
অঙ্গীকারনামা নথির যোগাযোগ অংশে রক্ষিত আছে,
দয়া করে দেখা যেতে পারে।
\\
\end{minipage}
\normalsize
\begin{minipage}[t]{0.05\linewidth}
% char
০৪।
\end{minipage}
\begin{minipage}[t]{0.95\linewidth}
\underline{\textbf{শুল্কায়ন সম্পর্কিত প্রস্তাব:}}
পণ্য চালান সংক্রান্ত আমদানি দলিলপত্র
ইনভয়েস পর্যালোচনা পূর্বক পণ্যের বর্ণনা, পরিমাণ, এইচ.এস.কোড, ঘোষিত মূল্য,
সমসাময়িককালের অভিন্ন/অনুরূপ পণ্যের একক মূল্যসহ নিম্নের ছকে শুল্কায়নের প্রস্তাব উপস্থাপন
করা হলো:
\end{minipage}
\scriptsize
{\taxtab}
\vspace*{1mm}
\normalsize
\noindent
\begin{minipage}[t]{0.05\linewidth}
% sat
০৫।
\end{minipage}
\begin{minipage}[t]{0.95\linewidth}
\underline{\textbf{শুল্কায়নযোগ্য মূল্য নিরুপন:}} ASYCUDA WORLD SYSTEM
সমসাময়িককালের অভিন্ন / অনুরূপ মূল্য তথ্য পর্যালোচনা করে দেখা
যায় যে, ইনভয়েস ঘোষিত মূল্যের চেয়ে ডাটাবেইস এর মূল্য উচ্চতর
বিধায় শুল্ক মূল্যায়ন বিধিমালা-২০০০ এর বিধি ৫ এর আওতায় প্রস্তাবিত
মূল্যে পণ্যচালানটি সাময়িক শুল্কায়ন করা যেতে পারে।
\\
\end{minipage}
\begin{minipage}[t]{0.05\linewidth}
% at
০৬।
\end{minipage}
\begin{minipage}[t]{0.95\linewidth}
\underline{\textbf{প্রস্তাব:}}
\end{minipage}
\footnotesize
\begin{minipage}[t]{0.05\linewidth}
\hspace{1em}
\end{minipage}
\begin{minipage}[t]{0.05\linewidth}
(ক)
\end{minipage}
\begin{minipage}[t]{0.90\linewidth}
{\srooot}, {\sroootd} মোতাবেক
ডি এফ ভ্যাট বাবদ ৪১,০৬১,.৭২ (একচল্লিশ হাজার একষট্টি টাকা বাহাত্তর পয়সা) টাকা
আদায় পূর্বক পণ্যচালানটি ঘোষিত HS Code
ও প্রস্তাবিত মূল্যে সাময়িক শুল্কায়ন করা যেতে
পারে।
\end{minipage}
\begin{minipage}[t]{0.05\linewidth}
\hspace{1em}
\end{minipage}
\begin{minipage}[t]{0.05\linewidth}
(খ)
\end{minipage}
\begin{minipage}[t]{0.90\linewidth}
কমিশনার, কাস্টমস মূল্যায়ন ও অভ্যন্তরীণ
নিরীক্ষা কমিশনারেট যথাসম্ভব স্বল্পতম সময়ে
নিরীক্ষা সমাপ্ত করিয়া নিরীক্ষা সমাপ্তির
১(এক) মাসের মধ্যে পরিশিষ্ট-৩ মোতাবেক একটি
বস্তুনিষ্ট প্রতিবেদন জাতীয় রাজস্ব বোর্ডের
সদস্য (কাস্টমস নীতি) এর নিকট এর অনুলিপি সংশ্লিষ্ট
কাস্টম হাউস বা কাস্টমস স্টেশন কর্তৃপক্ষ
বরাবরে প্রেরণ করিবেন।
নিরীক্ষায় রেয়াতী হারে আমদানিকৃত উপকরণের
ধর্তব্যযোগ্য অবব্যবহার পাওয়া গেলে
Customs Act এর Section-156
এর Sub Section (1) এর Table এর Item 10A
এর বিধান অনুযায়ী ব্যবস্থা গ্রহণ করিতে কাস্টমস
মূল্যায়ন ও অভ্যন্তরীন নিরীক্ষা কমিশনারেট আমদানি
সংশ্লিষ্ট কাস্টম কমিশনারকে অনুরোধ করিবেন।
প্রতিবেদন প্রাপ্তির পর যদি দেখা যায় যে, আমদানিকৃত
উপকরণ আমদানিকারকের প্রতিষ্ঠানে ব্যবহৃত হইয়াছে তাহা
হইলে অঙ্গীকারনামা ফেরতযোগ্য হইবে। উল্লেখ্য আবেদনকারী
উপকরণের প্রক্কলন দাখিল করা হয়েছে, যা নথির যোগাযোগ
অংশে রক্ষিত আছে।
\end{minipage}
\begin{minipage}[t]{0.05\linewidth}
\hspace{1em}
\end{minipage}
\begin{minipage}[t]{0.05\linewidth}
(গ)
\end{minipage}
\begin{minipage}[t]{0.90\linewidth}
বার্ষিক প্রাপ্যতার পরিমান
{\allotment} মে.টন।
ইতিপূর্বে ছাড় করা হয়েছে
{\relesed} মে.টন।
অবশিষ্ট আছে
({\allotment} - {\relesed}) মে.টন={\remaining} মে.টন।
বর্তমান চালানের পণ্য
{\current} মে.টন
ছাড় করা হলে
অবশিষ্ট থাকবে
({\remaining}-{\current}) মে.টন={\balance} মে.টন।
\end{minipage}
\begin{minipage}[t]{0.05\linewidth}
\hspace{1em}
\end{minipage}
\begin{minipage}[t]{0.05\linewidth}
(ঘ)
\end{minipage}
\begin{minipage}[t]{0.90\linewidth}
{\srooot}, {\sroootd} এর শর্ত মোতাবেক
শুল্ক মূল্যায়ন ও অভ্যন্তরীন নিরীক্ষা কমিশনারেট,
গুলফেশাঁ প্লাজা (৬ষ্ট তলা), ৬৯, আউটার
সার্কুলার রোড, মগবাজার, ঢাকা-১২১৭ এবং
আমদানিকারক বরাবর পত্র প্রেরণ করা যেতে পারে।
\\
\\
\normalsize
সদয় অবগতি ও আদেশার্থে উপস্থাপন করা
হলো।
\end{minipage}
\newpage
\begin{center}
\vspace*{6MM}
গণপ্রজাতন্ত্রী বাংলাদেশ সরকার
\\
\footnotesize{কাস্টম হাউস, চট্টগ্রাম।}
\\
\href{}{ওয়েব: www.chc.gov.bd}\hspace{1em}
\href{}{ইমেইল: customhousectg@gmail.com}
\end{center}
\normalsize
\begin{minipage}[t]{.74\linewidth}
{\fileno}
\end{minipage}
\begin{minipage}[t]{.26\linewidth}
{\fdt}
\\
\end{minipage}
\begin{minipage}[t]{.07\linewidth}
প্রাপক:
\end{minipage}
\begin{minipage}[t]{.93\linewidth}
কমিশনার
\\
শুল্ক মূল্যায়ন ও অভ্যন্তরীণ নিরীক্ষা কমিশনারেট
\\
গুলফেঁশা প্লাজা (৬ষ্ঠ তলা)
\\
৬৯, আউটার সার্কুলার রোড
\\
মগবাজার, ঢাকা-১২১৭।
\\
\end{minipage}
\begin{minipage}[t]{.07\linewidth}
বিষয়:
\end{minipage}
\begin{minipage}[t]{.93\linewidth}
\textbf{আমদানিকৃত পণ্য {\good} এর যথাযথ ব্যবহার নিশ্চিতকরণ।}
\end{minipage}
\begin{minipage}[t]{.07\linewidth}
সূত্র:
\end{minipage}
\begin{minipage}[t]{.93\linewidth}
{\srooot}, {\sroootd}।
\\
উপর্যুক্ত বিষয় ও সূত্রের প্রতি আপনার সদয় দৃষ্টি
আকর্ষন করা হলো।
\\
\end{minipage}
\begin{minipage}[t]{.07\linewidth}
২।
\end{minipage}
\begin{minipage}[t]{.93\linewidth}
আমদানিকারক প্রতিষ্ঠান- {\impn}, {\impaddsl}
কর্তৃক এলসি নং- {\lcno}, Date: {\lcdt}
এর মাধ্যমে {\co} হতে
{\pkg} সূত্রোক্ত জাহাজ যোগে {\good}
পণ্যচালানটি আমদানি করে পণ্য
খালাসের নিমিত্তে বি/ই নং- {\beno}, Date: {\bedt}
দাখিল করেন। আমদানিকারক একটি
শিল্প প্রতিষ্ঠান বিধায় পণ্যচালানটি
{\srooot}, {\sroootd}
এর শর্ত মোতাবেক আমদানি করেছেন এবং
তার প্রতিষ্ঠানটি বাংলাদেশ টেক্সটাইল মিলস
এসোসিয়েশন এর একটি সদস্য মর্মে
{\btmaltno}, {\btmaltnodt}
দাখিল করেছেন।
\\
\end{minipage}
\begin{minipage}[t]{.07\linewidth}
৩।
\end{minipage}
\begin{minipage}[t]{.93\linewidth}
উক্ত এসআরও এর শর্ত মোতাবেক
আলোচ্য আমদানিকারকের আমদানিকৃত পণ্য
শিল্প কারখানায় যথাযথভাবে ব্যবহার করা হয়েছে
কিনা তা আপনার দপ্তর কর্তৃক
যথাযথ স্বল্প সময়ে
নিরীক্ষা সমাপ্তির ৬(ছয়) মাসের মধ্যে এসআরও
এর পরিশিষ্ট-৩ মোতাবেক একটি বস্তুনিষ্ট
প্রতিবেদন জাতীয় রাজস্ব বোর্ডের সদস্য
(কাস্টমস নীতি) এর নিকট এবং অনুলিপি সংশ্লিষ্ট
কাস্টমস স্টেশন কর্তৃপক্ষ বরাবরে প্রেরণ করার
নির্দেশনা রয়েছে। উক্ত নির্দেশনা
অনুযায়ী প্রয়োজনীয় কর্যক্রম গ্রহণের জন্য
আপনাকে সবিনয় অনুরোধ করা হলো।
\\
\\
\\
\end{minipage}
\begin{minipage}[t]{0.60\linewidth}
\hspace{1em}
\end{minipage}
\begin{minipage}[t]{0.40\linewidth}
\begin{center}
সন্তোষ সরেন
\\
ডেপুটি কমিশনার অব কাস্টমস
\\
কমিশনার অব কাস্টমস এর পক্ষে
\\
কাস্টম হাউস, চট্টগ্রাম।
\\
\footnotesize{{\rodt}}
\vspace*{10MM}
\end{center}
\end{minipage}
\footnotesize{{\fileno}}
\\
\underline{\footnotesize{অনুলিপি অবগতি ও প্রয়োজনীয় ব্যবস্থা গ্রহনের জন্য প্রেরণ করা হলো}}
\\
\begin{minipage}[t]{0.06\linewidth}
\footnotesize{০১।}
\end{minipage}
\begin{minipage}[t]{0.94\linewidth}
\footnotesize{
আমদানিকারক {\impn}, {\impaddsl}
}
\end{minipage}
\begin{minipage}[t]{0.06\linewidth}
\footnotesize{০২।}
\end{minipage}
\begin{minipage}[t]{0.94\linewidth}
সভাপতি/সাধারণ সম্পাদক, বাংলাদেশ টেক্সটাইল
মিলস এসোসিয়েশন, ইউনিক ট্রেড সেন্টার (লেভেল-৮),
৮, পান্থপথ, কারওয়ান বাজার, ঢাকা-১২১৫।
\end{minipage}
\begin{minipage}[t]{0.06\linewidth}
\footnotesize{০৩।}
\end{minipage}
\begin{minipage}[t]{0.94\linewidth}
 অফিস কপি।
 \\
 \\
 \\
 \\
 \\
\end{minipage}
\begin{minipage}[t]{0.60\linewidth}
\hspace{1em}
\end{minipage}
\normalsize
\begin{minipage}[t]{0.40\linewidth}
\begin{center}
সন্তোষ সরেন
\\
ডেপুটি কমিশনার অব কাস্টমস
\\
কাস্টম হাউস, চট্টগ্রাম।
\end{center}
\end{minipage}
\thispagestyle{slogan}

\end{document}
